\section{Introduction}

Long-term behavioral studies of ants and other social insects that form large colonies present significant methodological challenges, particularly in quantifying individual abundance across extended observation periods. Traditional manual counting methods become prohibitively time-consuming and error-prone when dealing with densely populated colonies, especially in images or videos where individuals may overlap, be partially obscured, or be difficult to distinguish from complex backgrounds. These challenges are compounded in field studies where environmental factors such as variable lighting conditions, changing substrates, and weather-related disturbances can affect image quality and consistency over time.

Computer vision (CV) approaches offer promising solutions to automate individual counting and tracking, potentially enabling researchers to collect behavioral data at unprecedented temporal and spatial scales. However, successful implementation of CV systems in ecological research requires addressing several key practical considerations: determining minimum dataset requirements for model calibration and generalization, optimizing performance under challenging imaging conditions, and demonstrating the ecological insights that can be gained from automated data collection.

This study evaluates the application of CV techniques to automate ant counting in long-term behavioral studies, with particular emphasis on understanding foraging behavior dynamics over extended periods (weeks to months) in controlled environments. We address three critical objectives that collectively provide a framework for implementing CV-based counting systems in social insect research:

\begin{itemize}
    \item Objective 1: Determining the quantity of training data required for robust model generalization across varying imaging conditions, and investigating how background complexity affects model performance. This objective establishes practical guidelines for researchers designing imaging protocols for automated counting studies, addressing the common question of how much initial effort is needed to develop a functional system.
    \item Objective 2: Evaluating three different transfer learning approaches to adapt models trained on sparse colony images to densely packed scenarios where individuals are tightly clustered (e.g., up to one thousand ants per image). This objective demonstrates how researchers can efficiently extend their CV systems to new experimental conditions with minimal additional labeling effort, a critical consideration for long-term studies where colony density may change seasonally or experimentally.
    \item Objective 3: Demonstrating the application of automated counting results to understand spatial and temporal dynamics of ant foraging behavior through analysis of long-term image series. This objective showcases how CV-derived data can reveal ecological patterns that would be difficult to detect through traditional manual observation methods.
\end{itemize}

Although this study uses ant colonies as a model system, the methodological approach can be extended to other social species that form large aggregations, including bees, termites, wasps, and colonial vertebrates, provided researchers follow the systematic methodology outlined in Objectives 1 and 2 for dataset preparation and model optimization. The success of such applications depends primarily on establishing appropriate imaging protocols and following the data requirements and transfer learning strategies we demonstrate rather than on species-specific model architectures. To contextualize our approach within the broader field of ant behavioral ecology and to establish the foundation for our CV methodology, we first review the significance of foraging behavior in ant research, current methodological approaches to studying foraging dynamics, and the emerging applications of automated counting techniques in social insect ecology.

\section{Background study}

Ant foraging is one of the most studied themes in ant research because their behavioral repertoires are remarkably diverse \cite{reeves_evolution_2019}. Studying foraging behavior in ants can facilitate our understanding of ecological dynamics, as ants play vital roles in numerous ecosystem functions including nutrient cycling, seed dispersal, and habitat formation \cite{parr_response_2022, rocha_two_2024}. Their foraging patterns can also reveal insights into social organization, communication, and adaptability to environmental changes including pathogen challenge \cite{alciatore_immune_2021} or pesticide exposure \cite{thiel_sublethal_2016}. Foraging behaviors are particularly crucial when developing and evaluating ant baiting strategies for pest management, as the success of baiting relies heavily on ants ingesting a lethal dose of the bait, which can be facilitated by better understanding their foraging behavior \cite{galante_acute_2024}.

Quantifying long-term behaviors and investigating foraging preferences in ants typically involves manual counting across numerous images. This process is generally labor-intensive, time-consuming, and prone to human error, particularly when dealing with large numbers of ants. As a result, scaling up experiments to simulate more natural conditions, such as larger colony sizes and increased numbers of colonies, is challenging. For example, Hsu et al. \cite{hsu_viral_2018} observed foraging behavior on a total of 18 red imported fire ant (\emph{Solenopsis invicta}) colonies, with each colony comprising approximately 3,000 individual ants. Each colony was offered four food resources with different macronutrient ratios. The number of foragers on each food source was manually estimated by analyzing photos taken by an automatic recording system that captured images of each food resource every 10 minutes for 32 hours, resulting in a total of 13,824 photos processed by the study. Similarly, Kooij et al. \cite{kooij_differences_2014} manually observed ant foraging across multiple species, monitoring 9-11 trails per species at varying distances from the nest and at multiple times of the day. Although their method provided detailed insights into ant behavior and forage diversity, the manual nature of their approach potentially limited accuracy, repeatability, and scalability, underscoring the necessity of adopting automated counting solutions. Additionally, the number of ants counted in an image or at a specific time point may not fully capture the complexity of foraging behaviors, which often involve both spatial factors and the interaction between spatial and temporal dynamics, such as testing different sugar concentrations \cite{sola_feeding_2016} or bait types \cite{du_foraging_2023}.
    
These challenges make manual counting a laborious task. While it can still yield results, tools that specifically tackle this challenge would facilitate the counting process and thus the development of more effective pest control strategies. Relevant efforts have been made to automate the process of counting insects. One common approach, thresholding, is used to segment objects in images, enabling automated insect counting. This approach assumes that the object of interest exhibits distinct visual contrast with the background. For example, Parker et al. automatically counted black fly adults by placing them on Petri dishes with white backgrounds after collection from light traps \cite{parker_using_2020}. Similarly, Smythe et al. \cite{smythe_using_2020} used a comparable method to count horn flies on cattle in the field, relying on the color contrast between the flies and the cattle's skin. Software tools and computational pipelines have been released to facilitate the counting of birds \cite{marchowski_drones_2021} and tracking of moving animals \cite{chiara_animalta_2023} within designated areas. These methods typically rely on user-tuning approaches, such as adjusting image brightness or contrast to make objects of interest stand out from the background. While effective when examined images share consistent brightness or color contrast, these methods are sensitive to background variations and require careful manual parameter tuning to achieve optimal results each time a new image is processed. This limitation can be particularly challenging in long-term studies where backgrounds may change due to weather, lighting conditions, or camera repositioning.

A more advanced approach involves using deep learning models that learn the morphological patterns of objects and are less sensitive to background appearance. Such methods require standardized manual effort: labeling the location of each object of interest in a set of training images. Once trained, these models can automatically detect and count objects in new images without requiring manual parameter tuning. Additionally, the labeling efforts can be shared across different studies, providing broader utility compared to manually tuned parameters that are typically optimized for a specific dataset. For instance, the models YOLOv7 \cite{wang_yolov7_2022} and YOLOv8 \cite{jocher_yolo_2023,terven_comprehensive_2023} were leveraged to count polyphagous moths in customized traps \cite{saradopoulos_image-based_2023}. However, it can still be challenging to detect small objects like ants due to the limitations of current CV models, which are the primary approach for extracting information from images or videos. The challenges can be attributed to two main factors: pretraining data sources and model architecture. Typically, the pretraining data sources used to establish and train CV models are public datasets like COCO \cite{lin_microsoft_2015} and ImageNet \cite{deng_imagenet_2009}, designed for larger objects occupying a significant portion of an image. Since ants are small and often densely packed, standard models struggle to identify them accurately. Model architecture is another factor that limits the detection of small objects. Most CV models, such as YOLOv8, first resize input images to $640 \times 640$ pixels and then process them through several convolutional layers, resulting in progressively smaller feature maps. The smallest of these feature maps has a spatial dimension of $20 \times 20$ pixels. This downsampling means that an object smaller than $32 \times 32$ pixels (calculated by dividing 640 by 20) will be reduced to a single value in the smallest feature map. As a result, the model cannot effectively capture the fine spatial details necessary to distinguish small objects like ants, making accurate detection extremely difficult.
    
To address this limitation, a study introduced Fully Convolutional Regression Networks (FCRNs) to count small insects, such as thrips and aphids, on leaves \cite{bereciartua-perez_multiclass_2023}. The FCRN architecture is a symmetric encoder-decoder network that overcomes the size limitation of feature maps by upsampling them to the original image size, preserving the spatial information necessary for accurate counting. Additionally, the output of the FCRN is a density map that shares the same spatial dimensions as the input image and contains only the objects' centroids. Instead of predicting the exact size and location of each object, as in most standard object detection models, the FCRN estimates the probability of an object's presence at each pixel location. The disadvantage of this approach is that it usually requires a large amount of training data to generalize well to new images due to the lack of pretrained models on public datasets containing millions of images. Another strategy to improve the detection of small objects is slicing the input image into smaller patches and then feeding them separately into the model. This approach also overcomes the feature map size limitation, as it resizes the smaller patches to a larger size before processing \cite{hong_automatic_2021,bereciartua-perez_multiclass_2023}.